\section{The exact degree-ten flow closure}
\label{app:closure}

\paragraph{Targets and generators.}
At degrees four through ten the closure consumes exactly the generators
\begin{equation}
\operatorname{Tr}\tau,\;
\operatorname{Tr}\tau_2,\;
\operatorname{Tr}\tau_3,\;
\operatorname{Tr}\tau_4,\;
\operatorname{Tr}\tau_5,\;
(\operatorname{Tr}\tau)^2,\;
\operatorname{Tr}\tau\operatorname{Tr}\tau_2,\;
\operatorname{Tr}\tau\operatorname{Tr}\tau_3,\;
(\operatorname{Tr}\tau_2)^2,\;
\operatorname{Tr}\tau_2\operatorname{Tr}\tau_3 .
\end{equation}
This list was determined by instrumenting the closure rather than by reading the
construction, so that a generator consumed but not documented would show up.

\paragraph{Activation rules.}
A target is \emph{generated} when the flow produces it formally, and
\emph{activated} when the leading field degree of the generating expression is
low enough for it to contribute at the degree under consideration. The leading
degrees follow the rule: $\operatorname{Tr}(\tau)$ has leading field degree four
--- not two, by theorem~\ref{thm:gten} --- while $\operatorname{Tr}(\tau^k)$ has
$2k$ for $k \ge 2$, and products add. All generators were checked against this
rule independently at every prime, with no exceptions.

\paragraph{Why the distinction is the whole result.}
The span of \emph{generated} targets has dimension $\dimAten$, equal to
$\dim\Ften$. The span of \emph{activated} targets, closed under iteration, has
dimension $\dimDtenQ$. Identifying the two returns $\dim\Qten = 0$. The raw span
is retained in the test suite as a negative fixture for exactly this reason.

\paragraph{Fixed-point algorithm.}
Starting from the degree-four seed, repeatedly: generate targets from the current
span, admit those whose leading degree permits activation, and re-rank. Terminate
when a sweep raises no rank. Termination occurred after $\czSweeps$ sweeps.

\paragraph{The exact rational calculation.}
Of the closure's target rows, $\czIntegerRows$ were already integral. The
remaining $\czLiftedRows$ were reconstructed as follows: evaluate at
$\czFitPrimes$ fitting primes, combine by the Chinese remainder theorem, and
recover a rational by rational reconstruction~\cite{Wang:1982}. The reconstructed
rows were then validated at the held-out prime $\czHoldoutPrime$, which took no
part in the fit; there were no mismatches. The closure was then run entirely in
exact $\mathbb{Q}$ arithmetic.

\paragraph{Dimensions by degree.}
Over $\Q$ the reachable dimensions are $1$, $2$, $7$ and $\dimDtenQ$ at degrees
four, six, eight and ten.

\paragraph{Lower-bound certificate.}
An explicit $\czMinorSize \times \czMinorSize$ integer minor of the reachable
span has non-zero determinant. The submatrix, the column set, the integer scale
factor by which the rows were cleared of denominators, and the full integer
determinant are recorded in the artifact. The determinant was computed by Bareiss
elimination~\cite{Bareiss:1968} and confirmed by an independent exact rational
routine.

\paragraph{Upper bound: why none is needed.}
Running the closure over $\Q$ removes the modular admission test that made the
earlier figure a lower bound. At the fixed point, no further generated direction
raises the rank over $\Q$; therefore the reached span \emph{is} $\Dten$, and
$\dim_\Q\Dten = \dimDtenQ$ is an equality. This is the substantive difference
between the present computation and its modular predecessor.

\paragraph{Quotient maps.}
Three exact annihilating covectors are exhibited, each vanishing on $\Dten$ and
independent on $\Ften$. They give the quotient map to $\Qten$ directly, and the
quotient representatives are indexed by the free columns
(table~\ref{tab:quotient}).
